\chapter{线性表示、等价与不可约性}
\label{chap:representation-theory}

\chapref{chap:group-basics}讨论的是群本身的结构：群元怎样相乘、怎样分解、怎样通过正规子群与商群揭示内部组织。到了真正的物理应用里，我们还需要第二步：把抽象的群元变成线性空间上的可计算对象。点群要作用在晶格位移、原子轨道和振动模上，转动群要作用在波函数与态矢量上，置换群要作用在多体 Hilbert 空间上。于是我们自然会问：能否把一个抽象群 $G$ 送到某个线性空间上的可逆线性变换群中，并且保持群乘法？这正是群表示理论的出发点。

上一章的群元仍然是抽象对象。只知道 $gh=k$，我们能够研究群内部结构，却还不能把它直接放进量子力学方程或数值计算里。线性代数给了我们一个桥梁：把每个抽象群元 $g$ 对应到一个可逆线性变换 $\rho(g)$。这样，群乘法会变成矩阵乘法，抽象的“对称操作”就可以真正作用在坐标、轨道系数、振动位移或量子态上。

初学表示论时最容易混淆三个对象。第一层是抽象群元 $g\in G$；第二层是它在向量空间 $V$ 上对应的线性算符 $\rho(g):V\to V$；第三层是在选定一组基以后，这个线性算符写成的矩阵 $D(g)$。同一个 $\rho(g)$ 换一组基，矩阵会改变，但线性变换本身没有改变。表示论中所谓“两个表示等价”，很大程度上就是把这种无关紧要的坐标选择剔除掉。

本章最终要解决的问题可以用一句话概括：给定一个群在线性空间上的作用，怎样找出其中最小的、不能继续分裂的对称性块？这些最小块就是不可约表示。后面的能级简并、选择定则和能带标记之所以能用几个字母概括，本质上都是因为复杂的表示可以被拆成不可约表示。

除特别说明外，本章默认表示空间是有限维复线性空间 $V$，记 $\dim V=n$。这里“线性空间”意味着向量可以相加、可以乘复数；$\dim V=n$ 意味着存在 $n$ 个线性无关向量 $e_1,\dots,e_n$，使任意 $v\in V$ 都唯一写成
\[
v=v_1e_1+\cdots+v_ne_n.
\]
$\operatorname{GL}(V)$ 表示 $V$ 上全部\emph{可逆}线性变换组成的群；如果选定基，它与全部可逆 $n\times n$ 复矩阵组成的 $\operatorname{GL}(n,\CC)$ 对应。之所以使用复数域，是因为复数域既适合量子力学的波函数语言，又能保证有限群表示理论中若干关键结论以最简洁的形式成立。若选定 $V$ 的一组基，则线性变换可以写成矩阵；因此“群表示”也可理解为“把群元表示成矩阵并保持乘法”。但从逻辑上说，真正基本的是线性变换本身，矩阵只是基底选定后的坐标表达。
\section{表示的定义与最初例子}
\label{sec:representations-basic}


\begin{definition}[群表示]
设 $G$ 为群，$V$ 为有限维复线性空间。一个从 $G$ 到 $\operatorname{GL}(V)$ 的群同态
\[
\rho:G\longrightarrow \operatorname{GL}(V)
\]
称为 $G$ 在 $V$ 上的一个\textbf{线性表示}，简称\textbf{表示}。也就是说，对任意 $g,h\in G$，有
\[
\rho(gh)=\rho(g)\rho(h),
\qquad
\rho(e)=I_V.
\]
空间 $V$ 称为该表示的\textbf{表示空间}，其维数 $\dim V$ 称为表示的\textbf{维数}。
\label{def:representation}
\end{definition}

定义中的目标群必须是可逆线性变换群 $\operatorname{GL}(V)$，而不是所有线性变换的集合。原因与\chapref{chap:group-basics}中“群元必须有逆”完全一致：既然 $g^{-1}$ 在群里存在，$\rho(g)$ 也必须可逆，且其逆恰好是 $\rho(g^{-1})$。这一点直接由同态性质推出，也是后面许多运算的基础。

\begin{proposition}[表示的基本恒等式]
设 $\rho:G\to \operatorname{GL}(V)$ 是表示，则对任意 $g\in G$，有
\[
\rho(g^{-1})=\rho(g)^{-1}.
\]
若 $g^m=e$，则
\[
\rho(g)^m=I_V.
\]
特别地，有限群中每个群元都被表示成一个有限阶可逆线性变换。
\label{prop:basic-identities-rep}

\end{proposition}

\begin{proof}
由同态性质，
\[
\rho(g)\rho(g^{-1})=\rho(gg^{-1})=\rho(e)=I_V,
\]
同理也有 $\rho(g^{-1})\rho(g)=I_V$，故 $\rho(g^{-1})=\rho(g)^{-1}$。若 $g^m=e$，则反复应用同态性质即得
\[
\rho(g)^m=\rho(g^m)=\rho(e)=I_V.
\]
\end{proof}

一旦选定基底，表示就会转化为矩阵表示。这一步在实际计算中几乎总会做，但应当清楚：不同基底下得到的矩阵不同，它们描述的却是同一个表示。

\begin{definition}[矩阵表示]
设 $\rho:G\to\operatorname{GL}(V)$ 为表示，$\mathcal B=(e_1,\dots,e_n)$ 为 $V$ 的一组基。对每个 $g\in G$，记 $D_{\mathcal B}(g)$ 为线性变换 $\rho(g)$ 在基 $\mathcal B$ 下的矩阵，则有
\[
D_{\mathcal B}(gh)=D_{\mathcal B}(g)D_{\mathcal B}(h).
\]
映射 $g\mapsto D_{\mathcal B}(g)$ 称为表示 $\rho$ 在该基下的\textbf{矩阵表示}。
\label{def:matrix-representation}
\end{definition}

矩阵表示之所以重要，是因为之后的约化、对角化与特征标计算都发生在矩阵层面上；但如果过早把矩阵当成本体，就容易忽略基变换带来的冗余。因此本章反复区分“表示”与“某一基下的表示矩阵”。

\begin{definition}[子表示与不变子空间]
设 $\rho:G\to\operatorname{GL}(V)$ 为表示，$W\subseteq V$ 为线性子空间。若对任意 $g\in G$ 都有
\[
\rho(g)W\subseteq W,
\]
则称 $W$ 为 $\rho$ 的\textbf{不变子空间}。这时限制映射
\[
\rho\vert_W:G\to\operatorname{GL}(W)
\]
给出一个新的表示，称为原表示的\textbf{子表示}。
\label{def:subrepresentation}
\end{definition}

不变子空间的概念看似简单，却是整章最核心的几何对象。因为表示是否"能被分解"，实质上就是问：除了 $\{0\}$ 与 $V$ 本身之外，它有没有非平凡不变子空间？这个问题将在下面变成"可约与不可约"的定义。

在具体问题里，表示主要来自三种渠道：群本身已经由矩阵组成；群通过变换某个集合而诱导出向量空间上的作用；以及物理对象本身带有某个天然的线性空间结构。先看最基本的几个例子。

\medskip
\noindent\emph{平凡表示。}
对任意群 $G$，在一维空间 $\mathbb C$ 上定义
\[
\rho_{\mathrm{triv}}(g)=1,
\qquad \forall g\in G.
\]
这给出一个一维表示，称为\textbf{平凡表示}。它说明并不是每个表示都携带复杂结构，但它在后面的特征标分解中会反复出现。
\par\medskip

\medskip
\noindent\emph{由群作用诱导的置换表示。}
设群 $G$ 作用在有限集合 $X=\{x_1,\dots,x_n\}$ 上。把以 $X$ 为基的向量空间记为
\[
V_X=\operatorname{span}_{\mathbb C}\{e_{x_1},\dots,e_{x_n}\}.
\]
对每个 $g\in G$ 定义线性变换
\[
\rho_X(g)e_x=e_{g\cdot x}.
\]
由于群作用满足 $(gh)\cdot x=g\cdot(h\cdot x)$ 与 $e\cdot x=x$，所以 $\rho_X$ 是一个表示，称为由该群作用诱导出的\textbf{置换表示}。这个例子直接把\chapref{chap:group-basics}的群作用与本章的表示联系起来：群先作用在集合上，再线性延拓到以该集合为基的向量空间上。
\label{ex:permutation-representation}
\par\medskip

\begin{example}[$S_3$ 的二维标准表示]
上一章我们已经知道 $S_3$ 是三阶置换群。它对三维空间 $\mathbb C^3$ 的自然作用是置换标准基向量 $e_1,e_2,e_3$，即得到三维置换表示。这个表示中，
\[
u_0=e_1+e_2+e_3
\]
张成的一维子空间在所有置换下保持不变，因此给出一个平凡子表示。与之互补的二维子空间
\[
W=\{(z_1,z_2,z_3)\in\mathbb C^3\mid z_1+z_2+z_3=0\}
\]
也在 $S_3$ 作用下不变，故限制到 $W$ 上得到一个二维表示。这个表示就是 $S_3$ 最重要的不可约表示之一，后面构造其特征标表时还会再次遇到。这个例子说明：我们之所以关心不变子空间，是因为它能把一个较大的表示分裂成更小、更本质的成分。
\end{example}

为了帮助读者把“抽象群 $\to$ 线性变换群”这层关系视觉化，\figref{fig:representation-homomorphism} 给出了一个最简示意。它并不添加新的数学内容，但有助于在后文反复出现基变换、交织算子与特征标时保持概念清晰。

\begin{figure}[htbp]
\centering
\begin{tikzfig}
  \node[lecturebox] (G) at (0,0) {$G$\\抽象群};
  \node[lecturebox] (GL) at (7.0,0) {$\operatorname{GL}(V)$\\线性变换群};
  \draw[lecturearrow] (G) -- node[above,lecturelabel] {$\rho$} (GL);
  \node[lecturelabel] at (3.5,-1.2) {$\rho(gh)=\rho(g)\rho(h)$,\quad $\rho(e)=I_V$};
  \node[lecturebox] (M) at (7.0,-3.0) {固定基后的矩阵群\\$D(g)$};
  \draw[lecturearrow] (GL) -- node[right,lecturelabel] {选定基} (M);
\end{tikzfig}
\caption{表示的逻辑层次：先有表示 $\rho:G\to \operatorname{GL}(V)$，再在选定基后得到矩阵表示。}
\label{fig:representation-homomorphism}
\end{figure}

把“抽象群元 $\to$ 线性算符 $\to$ 矩阵”三层关系放到 $C_3$ 上看最清楚。取生成元 $a$，满足 $a^3=e$。在二维实平面上，让 $a$ 表示逆时针转 $120^\circ$，则在标准基下
\[
D(a)=
\begin{pmatrix}
-\tfrac12&-\tfrac{\sqrt3}{2}\\
\tfrac{\sqrt3}{2}&-\tfrac12
\end{pmatrix},
\qquad
D(a)^3=I_2.
\]
于是 $D(a^2)=D(a)^2$、$D(e)=I_2$，整个表示只要知道生成元矩阵就确定了。

如果改到复数基
\[
f_+=\frac1{\sqrt2}(e_x-\ii e_y),\qquad
f_-=\frac1{\sqrt2}(e_x+\ii e_y),
\]
同一个转动会变成对角矩阵
\[
D'(a)=
\begin{pmatrix}
\ee^{2\pi\ii/3}&0\\
0&\ee^{-2\pi\ii/3}
\end{pmatrix}.
\]
这两个矩阵看起来完全不同，却由同一个基变换联系，描述同一个表示。更重要的是，复基已经把二维表示拆成两个一维不变子空间。这预示了“基变换、等价表示、不可约分解”其实是同一个问题的不同层面。

\section{等价表示、不可约表示与酉表示}
\label{sec:equiv-irreducible-unitary}


假设同一个二维几何变换，在一组直角坐标基下写成矩阵 $D(g)$，换到另一组斜基以后写成 $D'(g)$。矩阵的数字当然不同，但我们不会说物理对称性变了。于是表示论首先要建立一个“忽略坐标选择”的等价关系。

若两个表示只是在不同基下写出的同一个线性作用，那么把它们当成不同对象会造成大量重复。于是首先要消除“基的任意性”带来的冗余，这就得到等价表示的概念。

\begin{definition}[交织算子与表示等价]
设 $(\rho,V)$ 与 $(\rho',V')$ 是群 $G$ 的两个表示。若线性映射
\[
T:V\to V'
\]
满足对任意 $g\in G$ 都有
\[
T\rho(g)=\rho'(g)T,
\]
则称 $T$ 是这两个表示之间的\textbf{交织算子}。若存在可逆交织算子 $T$，则称这两个表示\textbf{等价}，记作
\[
(\rho,V)\simeq(\rho',V').
\]
\label{def:equivalent-representations}
\end{definition}

这个定义看似抽象，其实完全是在形式化“同一个表示在不同基下有不同矩阵”的直觉。交织关系要求先作用群元再做 $T$，与先做 $T$ 再作用群元的结果相同；而 $T$ 可逆则保证两个表示的信息完全一致。

\begin{proposition}[基变换公式]
设 $\rho:G\to\operatorname{GL}(V)$ 是表示，$\mathcal B$ 与 $\mathcal B'$ 是 $V$ 的两组基，$S$ 为从坐标向量 $[v]_{\mathcal B'}$ 到 $[v]_{\mathcal B}$ 的变换矩阵，即
\[
[v]_{\mathcal B}=S[v]_{\mathcal B'}.
\]
则对应的矩阵表示满足
\[
D_{\mathcal B'}(g)=S^{-1}D_{\mathcal B}(g)S.
\]
因此同一表示在不同基下总给出等价的矩阵表示。
\label{prop:change-of-basis}

\end{proposition}

\begin{proof}
由定义，若 $v\in V$，则
\[
[\rho(g)v]_{\mathcal B}=D_{\mathcal B}(g)[v]_{\mathcal B},
\qquad
[\rho(g)v]_{\mathcal B'}=D_{\mathcal B'}(g)[v]_{\mathcal B'}.
\]
又有
\[
[\rho(g)v]_{\mathcal B}=S[\rho(g)v]_{\mathcal B'}
\]
以及 $[v]_{\mathcal B}=S[v]_{\mathcal B'}$，故
\[
S D_{\mathcal B'}(g)[v]_{\mathcal B'}
= D_{\mathcal B}(g)S[v]_{\mathcal B'}
\]
对任意坐标向量成立，于是
\[
SD_{\mathcal B'}(g)=D_{\mathcal B}(g)S.
\]
两边左乘 $S^{-1}$，得到 $D_{\mathcal B'}(g)=S^{-1}D_{\mathcal B}(g)S$，即得所求。
\end{proof}

表示等价的交换图如图 \figref{fig:intertwiner-diagram} 所示。
\begin{figure}[htbp]
\centering
\begin{tikzcd}[column sep=large,row sep=large]
V \arrow[r,"\rho(g)"] \arrow[d,"T"'] & V \arrow[d,"T"] \\
V' \arrow[r,"\rho'(g)"'] & V'
\end{tikzcd}
\caption{交织算子满足的交换图。若 $T$ 可逆，则两个表示等价。}
\label{fig:intertwiner-diagram}
\end{figure}

有了等价表示的语言，下一个核心问题是：什么样的表示"不能被拆开"？这就是不可约表示。


为什么要找不变子空间？设一组表示矩阵在某个基下可以同时写成
\[
D(g)=\begin{pmatrix}D_1(g)&0\\0&D_2(g)\end{pmatrix}
\qquad \forall g\in G.
\]
那么前一组坐标和后一组坐标在所有群操作下永远不会互相混合。也就是说，一个本来较大的对称性问题实际上分成两个互不耦合的小问题。继续这样分，直到再也分不下去，所得块就是不可约表示。

\chapref{chap:group-basics}中我们把一个群通过子群、正规子群与商群来分解；在表示论里，对应的分解工具是不变子空间。于是“是否还能继续分解”被正式表述为可约与不可约。

\begin{definition}[可约表示与不可约表示]
设 $(\rho,V)$ 为群 $G$ 的表示。若存在非平凡不变子空间
\[
\{0\}\subsetneq W\subsetneq V,
\]
则称 $\rho$ 为\textbf{可约表示}。若不存在这样的 $W$，则称 $\rho$ 为\textbf{不可约表示}（不可约表示，常缩写为 \textbf{irrep}）。
\label{def:irreducible}
\end{definition}

不可约表示之于表示论，正如素数之于整数分解。一般表示的真正结构并不体现在“矩阵有多大”，而体现在它能否分裂为若干不可约成分，以及这些成分各自出现多少次。为了说明这一点，先记一个最基本的直接和结构。

\begin{definition}[表示的直接和]
设 $(\rho_1,V_1)$ 与 $(\rho_2,V_2)$ 是群 $G$ 的表示。在直和空间 $V_1\oplus V_2$ 上定义
\[
(\rho_1\oplus \rho_2)(g)(v_1,v_2)=\bigl(\rho_1(g)v_1,\rho_2(g)v_2\bigr).
\]
这给出新的表示，称为 $\rho_1$ 与 $\rho_2$ 的\textbf{直接和}。
\label{def:direct-sum-representation}
\end{definition}

若在某组基下 $V_1$ 与 $V_2$ 对应于前后两个坐标块，则矩阵形式恰好是分块对角矩阵
\[
D_{\rho_1\oplus\rho_2}(g)=
\begin{pmatrix}
D_{\rho_1}(g)&0\\
0&D_{\rho_2}(g)
\end{pmatrix}.
\]
因此，可约性在矩阵上就是“能否通过同一个基变换把所有表示矩阵同时化为分块对角”的问题。

\medskip
\noindent\emph{$S_3$ 的三维置换表示如何分解。}
前面已经指出，$S_3$ 在 $\mathbb C^3$ 上的置换表示存在一维不变子空间 $\mathbb C u_0$，其中 $u_0=(1,1,1)$。如果取
\[
W=\{(z_1,z_2,z_3)\in\mathbb C^3\mid z_1+z_2+z_3=0\},
\]
则 $\mathbb C^3=\mathbb C u_0\oplus W$，且这两个子空间都不变。因此三维置换表示分解为
\[
\rho_{\mathrm{perm}}\simeq \rho_{\mathrm{triv}}\oplus \rho_{\mathrm{std}},
\]
其中 $\rho_{\mathrm{std}}$ 是二维标准表示。这个例子说明：所谓“找到不变子空间”，最终目的并不是停留在几何描述，而是把表示拆成更小的块。
\par\medskip

可约性也可以直接用矩阵读出来。假设对每个 $g$ 都有
\[
D(g)=
\begin{pmatrix}
* & * & 0\\
* & * & 0\\
0 & 0 & *
\end{pmatrix}.
\]
那么前两个基向量张成的二维空间在群作用下不会跑到第三个方向，而第三个基向量张成的一维空间也保持自己，所以表示至少分成 $2+1$ 两块。相反，如果只是某一个 $D(g)$ 可以对角化，并不能说明表示可约；关键要求是\emph{同一个子空间对所有 $g\in G$ 都不变}，也就是所有表示矩阵能在同一组基下同时形成共同的块结构。

在实际应用中，我们几乎总是与酉表示打交道——量子力学中的对称变换是酉的，有限群的表示总可以酉化。酉性不仅保证内积不变，还使可约表示自动分解为正交直和，是 Maschke 定理的几何核心。

在线性代数里，内积 $\langle u,v\rangle$ 同时编码长度与夹角；在量子力学里，它进一步决定概率幅。一个线性变换若保持内积，就不会改变向量长度，也不会改变任意两向量之间的夹角，这样的变换称为酉变换。因此“对称操作应保持概率”自然把量子力学中的表示引向酉表示。

在量子力学里，保持内积往往比仅仅可逆更重要，因此我们希望把表示尽量放到酉变换的框架内。一般群并不保证每个有限维表示都天然酉，但对有限群却总可以通过平均化把内积改造成群不变的，于是任何有限群表示都等价于酉表示。

\begin{definition}[酉表示]
设 $V$ 为带厄米内积 $\langle\cdot,\cdot\rangle$ 的复内积空间。若表示 $\rho:G\to\operatorname{GL}(V)$ 满足对任意 $g\in G$ 与 $u,v\in V$ 有
\[
\langle \rho(g)u,\rho(g)v\rangle=\langle u,v\rangle,
\]
则称 $\rho$ 为\textbf{酉表示}。等价地，每个 $\rho(g)$ 都是酉算子，即
\[
\rho(g)^{\dagger}=\rho(g)^{-1}.
\]
\label{def:unitary-representation}
\end{definition}

\begin{theorem}[有限群表示的酉化]
设 $G$ 为有限群，$(\rho,V)$ 为其有限维复表示。则存在 $V$ 上的某个厄米内积，使得 $\rho$ 成为酉表示。特别地，任何有限群的有限维复表示都等价于酉表示。
\label{thm:unitarization}

\end{theorem}

\begin{proof}
先任取 $V$ 上一个正定厄米内积 $(\cdot,\cdot)_0$。利用有限群上的平均，定义新内积
\[
\langle u,v\rangle
=\frac1{|G|}\sum_{g\in G}(\rho(g)u,\rho(g)v)_0.
\]
线性与共轭线性由原内积逐项继承。由于每一项都非负，且当 $u\neq 0$ 时取 $g=e$ 的那一项严格正，因此该内积仍正定。

现在验证它对群作用不变。对任意固定的 $h\in G$，有
\begin{align*}
\langle \rho(h)u,\rho(h)v\rangle
&=\frac1{|G|}\sum_{g\in G}(\rho(g)\rho(h)u,\rho(g)\rho(h)v)_0\\
&=\frac1{|G|}\sum_{g\in G}(\rho(gh)u,\rho(gh)v)_0.
\end{align*}
当 $g$ 走遍 $G$ 时，$gh$ 也恰好走遍 $G$ 中全部元素一次，因此把求和指标换成 $k=gh$，得
\[
\langle \rho(h)u,\rho(h)v\rangle
=\frac1{|G|}\sum_{k\in G}(\rho(k)u,\rho(k)v)_0
=\langle u,v\rangle.
\]
故 $\rho$ 在新内积下是酉表示。
\end{proof}

平均化方法可以再拆开看一次。任意初始内积 $(\cdot,\cdot)_0$ 通常不满足群不变性，也就是 $(\rho(h)u,\rho(h)v)_0$ 可能依赖于 $h$。我们把所有群元产生的这些内积全部平均：
\[
\langle u,v\rangle_G=\frac1{|G|}\sum_{g\in G}(\rho(g)u,\rho(g)v)_0.
\]
现在再作用一个固定的 $h$：
\[
\langle\rho(h)u,\rho(h)v\rangle_G
=\frac1{|G|}\sum_{g\in G}(\rho(gh)u,\rho(gh)v)_0.
\]
因为右乘 $h$ 只是把有限集合 $G$ 的元素重新排列一次，$gh$ 仍然遍历全部群元，所以这个求和与原求和完全相同。这就是群平均能够制造不变量的核心原因。

以后看到
\[
\frac1{|G|}\sum_{g\in G}(\cdots)
\]
时，应立刻联想到“把一个任意对象平均成群不变对象”。投影算符、大正交关系和许多对称化构造都在重复这同一个思想。


这个平均化证明非常值得反复体会。它与\chapref{chap:group-basics}中从群平均构造投影、以及后面 Maschke 定理中的平均化方法是同一个思想：如果群是有限的，就可以把任意选择的对象“平均”成群不变对象。对表示论而言，这一技巧几乎贯穿始终。

\begin{proposition}[酉表示的正交补不变]
设 $(\rho,V)$ 为酉表示，$W\subseteq V$ 为不变子空间，则其正交补
\[
W^{\perp}=\{v\in V\mid \langle v,w\rangle=0,\ \forall w\in W\}
\]
也是不变子空间。
\label{prop:orthogonal-complement-invariant}

\end{proposition}

\begin{proof}
取 $v\in W^{\perp}$、$w\in W$、$g\in G$。因为 $W$ 不变，所以 $\rho(g^{-1})w\in W$。又因表示酉，故
\[
\langle \rho(g)v,w\rangle
=\langle v,\rho(g)^{\dagger}w\rangle
=\langle v,\rho(g^{-1})w\rangle=0.
\]
于是 $\rho(g)v\in W^{\perp}$，故 $W^{\perp}$ 不变。
\end{proof}

这一定理说明：在酉表示里，一旦找到一个不变子空间，就自动得到其不变互补，从而能够真正把表示拆成直接和。这正是有限群表示完全可约的核心几何原因；Maschke 定理会把它推广成一个完整的结构定理。

不可约表示是表示论的"原子"，而 Schur 引理则是判断不可约性的基本工具：它断言，不可约表示之间只有"平凡"的交织算符。

现在已经知道不可约表示是“再也拆不开”的最小块。接下来问一个更细的问题：如果一个线性算符 $T$ 与群中每个表示矩阵都交换，即
\[
T\rho(g)=\rho(g)T\qquad\forall g\in G,
\]
那么 $T$ 在一个不可约块中还能有多复杂？直觉上，如果 $T$ 有不同本征空间，而这些本征空间又被群作用保持，那么不可约性就会被破坏。Schur 引理把这个直觉变成严格结论：在复数域上的不可约表示里，这样的 $T$ 只能是恒等算符的标量倍数。

不可约表示的内部同态极其刚性，这一事实由 Schur 引理给出。它既是理论结构的核心，也是推导大正交关系与特征标正交关系的关键工具。

\begin{theorem}[Schur 引理]
设 $(\rho,V)$ 与 $(\sigma,W)$ 是群 $G$ 的两个有限维不可约复表示。
\begin{enumerate}[label=(\roman*),leftmargin=3em]
  \item 若 $T:V\to W$ 是交织算子，则 $\operatorname{Ker} T$ 与 $\operatorname{Im} T$ 分别是 $V$ 与 $W$ 的不变子空间。因此，$T$ 要么为零算子，要么为同构。
  \item 若进一步 $V=W$ 且 $\rho=\sigma$，则任一交织算子 $T:V\to V$ 必为某个标量的倍数：
  \[
  T=\lambda I_V.
  \]
\end{enumerate}
\label{thm:schur-lemma}

\end{theorem}

\begin{proof}
(i) 先证核不变。若 $v\in \operatorname{Ker} T$，则
\[
T\rho(g)v=\sigma(g)Tv=0,
\]
故 $\rho(g)v\in\operatorname{Ker} T$。再证像不变。若 $w\in\operatorname{Im} T$，写成 $w=Tv$，则
\[
\sigma(g)w=\sigma(g)Tv=T\rho(g)v\in\operatorname{Im} T.
\]
由于 $V,W$ 都不可约，$\operatorname{Ker} T$ 只能是 $\{0\}$ 或 $V$，$\operatorname{Im} T$ 只能是 $\{0\}$ 或 $W$。若 $T\neq 0$，则 $\operatorname{Ker} T\neq V$，只能是 $\{0\}$，同时 $\operatorname{Im} T\neq\{0\}$，只能是 $W$，于是 $T$ 为双射。

(ii) 取 $T$ 的任一本征值 $\lambda$（复数域上总存在），考虑算子
\[
S=T-\lambda I_V.
\]
由于 $T$ 与 $I_V$ 都交织，故 $S$ 也交织。又因为存在本征向量 $v\neq 0$ 满足 $Sv=0$，故 $S$ 不是同构。由 (i) 得 $S=0$，因此 $T=\lambda I_V$。
\end{proof}

Schur 引理的第一部分其实只用了一个非常简单的观察：交织算子的核和像会被群作用保持。若表示不可约，那么不变子空间只允许 $\{0\}$ 与整个空间，因此一个非零交织算子既不能有非平凡核，也不能有非满的像，只能是同构。

第二部分为什么需要复数域？因为我们要从 $T$ 中取一个本征值 $\lambda$。在有限维复向量空间上，特征多项式一定有根，所以至少存在一个非零向量 $v$ 满足 $Tv=\lambda v$。于是 $T-\lambda I$ 有非零核，不可能是同构；第一部分立刻迫使它为零，从而 $T=\lambda I$。如果底域是实数，这一步可能没有实本征值，所以结论需要修改。


Schur 引理告诉我们：不可约表示一旦固定，其内部与群作用相容的线性算子少得惊人，几乎只有标量。这就是为什么很多“平均得到的交织算子”最后都会塌缩成某个常数乘恒等算子；后面正交关系中的归一化常数正是这样算出来的。

Schur 引理在量子力学中的一个直接原型是：若一个不可约多重态上有算符 $A$ 与所有对称操作都交换，那么 $A$ 在这个多重态内只能是常数倍单位矩阵。若 $A$ 就是哈密顿量，这意味着一个 $d$ 维不可约表示中的 $d$ 个态必须具有同一能量。第四章会把这句话严格写成能级简并定理。
